\documentclass[11pt]{article}

\usepackage[margin=1in]{geometry}
\usepackage{amsmath, amssymb}
\usepackage{bm}
\usepackage{booktabs}
\usepackage{hyperref}
\usepackage{enumitem}

\newcommand{\R}{\mathbb{R}}
\newcommand{\quat}[1]{\bm{#1}}
\newcommand{\vp}{\bm{p}}
\newcommand{\vv}{\bm{v}}
\newcommand{\vw}{\bm{\omega}}
\DeclareMathOperator{\rotvec}{rotvec}
\DeclareMathOperator{\diag}{diag}

\title{\texttt{quest\_teleop}: Control Logic and Math}
\author{UR5 + Robotiq Quest~3S teleoperation package}
\date{\today}

\begin{document}
\maketitle

\begin{abstract}
This document specifies the control math implemented in the
\texttt{quest\_teleop} package (\texttt{src/quest\_teleop/}). The mapping from
Meta Quest~3S controller poses to UR5 motion follows the DROID
\texttt{VRPolicy} formulation. One backend-agnostic control loop
(\texttt{teleop.py}) drives two interchangeable robot backends: a direct
\texttt{ur\_rtde} backend (real robot / URSim) and a ROS~2 backend
(\texttt{/forward\_velocity\_controller}, for the \texttt{launch\_all.sh}
fake-hardware simulation). The two backends share \emph{identical} control
math; they differ only in how the final Cartesian command is realised
(Section~\ref{sec:backends}).
\end{abstract}

\tableofcontents

\section{Notation and frames}
Quaternions are scalar-first, $\quat{q}=(w,x,y,z)$, unit norm, acting as
rotations via the Hamilton product $\otimes$. The conjugate is
$\quat{q}^{*}=(w,-x,-y,-z)$. The map $\rotvec(\quat{q})\in\R^3$ returns the
axis--angle (rotation) vector and $\rotvec^{-1}$ its inverse
(\texttt{quat\_to\_rotvec}, \texttt{rotvec\_to\_quat}).

\begin{description}[leftmargin=2.4em,style=nextline]
  \item[VR raw frame] $H^{\text{raw}}\in SE(3)$: the controller pose as
        reported by \texttt{oculus\_reader} (a $4\times4$ homogeneous matrix).
  \item[Global$\to$env remap] $M\in\R^{4\times4}$: a signed axis permutation
        (\texttt{vec\_to\_reorder\_mat}) aligning the VR frame to the robot
        base frame. Default reorder $[-2,-1,-3,4]$.
  \item[Alignment] $V\in SE(3)$: the ``forward reset'' transform
        (Section~\ref{sec:yaw}), a pure rotation.
  \item[Robot base frame] all robot poses (TCP, origins, targets) live here.
\end{description}

The controller pose expressed in the robot/env frame is
\begin{equation}
  H = M \, V \, H^{\text{raw}},\qquad
  \vp_{\text{vr}} = H_{1:3,4},\quad
  \quat{q}_{\text{vr}} = \mathrm{rmat2quat}(H_{1:3,1:3}).
  \label{eq:process}
\end{equation}

\section{Forward reset: yaw-only alignment}\label{sec:yaw}
A joystick click (or the initial orientation reset) sets the alignment $V$ so
the controller's heading maps to the robot's forward direction. To keep
gravity vertical regardless of how the controller is tilted at click time, we
keep \emph{only} the yaw (rotation about the env-vertical axis) and discard the
tilt --- a swing--twist decomposition.

Let the env ``up'' axis be $\bm{e}_z=(0,0,1)$ in the robot frame; the
corresponding axis in the VR frame is
$\bm{a} = M_{1:3,1:3}^{\top}\,\bm{e}_z$. Given the click-time controller
rotation quaternion $\quat{q}=(w,\bm{u})$ with vector part $\bm{u}$, the twist
about $\bm{a}$ is
\begin{equation}
  \quat{q}_{\text{yaw}} = \frac{(w,\;(\bm{u}\cdot\hat{\bm a})\,\hat{\bm a})}
       {\lVert (w,\;(\bm{u}\cdot\hat{\bm a})\,\hat{\bm a})\rVert},
  \qquad \hat{\bm a}=\bm{a}/\lVert\bm{a}\rVert
  \label{eq:twist}
\end{equation}
(\texttt{twist\_about\_axis}). The alignment is the inverse yaw rotation,
$V = \diag\!\big(R(\quat{q}_{\text{yaw}}^{*}),\,1\big)$, with zero translation;
the translational baseline is handled separately by the origins
(Section~\ref{sec:droid}).

\section{VR pose low-pass filter}\label{sec:filter}
A one-pole filter smooths the controller pose before it drives the robot
(\texttt{filters.VRPoseFilter}, ported from Open-Teach). With smoothing
constant $\alpha\in[0,1)$ (\texttt{--filter-alpha}, default $0.8$),
\begin{align}
  \vp \;&\leftarrow\; \alpha\,\vp + (1-\alpha)\,\vp_{\text{vr}}, \\
  \quat{q} \;&\leftarrow\; \mathrm{slerp}\!\big(\quat{q},\,\quat{q}_{\text{vr}},\,1-\alpha\big).
\end{align}
$\alpha=0$ disables filtering; larger $\alpha$ is smoother but laggier. The
first sample after a re-baseline passes through unfiltered so the origins
below are captured from the same (filtered) stream.

\section{DROID target pose}\label{sec:droid}
While the grip deadman is held, origins are latched on the first tick:
the robot TCP origin $(\vp_{o}^{R},\quat{q}_{o}^{R})$ and the (filtered) VR
origin $(\vp_{o}^{V},\quat{q}_{o}^{V})$. The commanded target tracks the
controller's \emph{displacement} from its origin.

\paragraph{Position (with scaling).}
\begin{equation}
  \vp_{\text{tgt}} = \vp_{o}^{R} + s\,\big(\vp - \vp_{o}^{V}\big),
  \label{eq:tgtpos}
\end{equation}
where $s = s_{\text{pos}}\cdot(s_{\text{prec}}\ \text{if precision else }1)$
(\texttt{--pos-scale}, default $1.0$; precision toggled by the A/X button,
\texttt{--precision-scale}, default $0.5$).

\paragraph{Orientation.} The controller's incremental rotation is applied onto
the robot origin orientation:
\begin{equation}
  \Delta\quat{q} = \quat{q} \otimes (\quat{q}_{o}^{V})^{*},\qquad
  \quat{q}_{\text{tgt}} = \Delta\quat{q} \otimes \quat{q}_{o}^{R}.
  \label{eq:tgtquat}
\end{equation}

Equations \eqref{eq:tgtpos}--\eqref{eq:tgtquat} are shared by both control
modes and both backends.

\section{Control modes}
Let the current TCP pose be $(\vp_{c},\quat{q}_{c})$ and the control period
$\Delta t = 1/f$ (\texttt{--hz}, default $50$).

\subsection{Speed mode (\texttt{--mode speed}, default)}
A proportional controller on the pose error produces a Cartesian twist:
\begin{align}
  \vv &= k_p\,(\vp_{\text{tgt}} - \vp_{c}), &
  k_p &= \texttt{pos\_gain}\ (3.0),\\
  \quat{e} &= \quat{q}_{\text{tgt}}\otimes\quat{q}_{c}^{*}\ \text{(shortest)}, &
  \vw &= k_r\,\rotvec(\quat{e}), \quad k_r=\texttt{rot\_gain}\ (2.0).
\end{align}

\paragraph{Uniform saturation.} When either component exceeds its limit, both
are scaled by a single factor so the screw axis (motion direction) is
preserved (\texttt{limit\_velocity}, after DROID):
\begin{equation}
  \sigma = \min\!\left(1,\ \frac{v_{\max}}{\lVert\vv\rVert},\
                       \frac{\omega_{\max}}{\lVert\vw\rVert}\right),\qquad
  \vv\leftarrow\sigma\vv,\quad \vw\leftarrow\sigma\vw,
\end{equation}
with $v_{\max}=\texttt{max\_lin\_vel}\ (0.15\,\mathrm{m/s})$,
$\omega_{\max}=\texttt{max\_rot\_vel}\ (0.75\,\mathrm{rad/s})$. The twist
$(\vv,\vw)$ is the command handed to the backend.

\subsection{Servo mode (\texttt{--mode servo}, RTDE only)}
Instead of a velocity, a clamped absolute target pose is streamed. A per-cycle
displacement clamp (after the Open-Teach/deoxys operator) bounds the step:
\begin{align}
  \vp_{\text{cmd}} &= \vp_{c} + \mathrm{clip}\big(\vp_{\text{tgt}}-\vp_{c},\,\delta_p\big),\\
  \vw_{\text{step}} &= \mathrm{clip}\big(\rotvec(\quat{q}_{\text{tgt}}\otimes\quat{q}_{c}^{*}),\,\delta_r\big),\quad
  \quat{q}_{\text{cmd}} = \rotvec^{-1}(\vw_{\text{step}})\otimes\quat{q}_{c},
\end{align}
where $\mathrm{clip}(\bm{x},\delta)$ scales $\bm{x}$ to norm $\le\delta$
(\texttt{scale\_to\_max\_norm}). Defaults
$\delta_p=v_{\max}\Delta t$, $\delta_r=\omega_{\max}\Delta t$
(\texttt{--max-pos-step}, \texttt{--max-rot-step}). The pose
$(\vp_{\text{cmd}},\rotvec(\quat{q}_{\text{cmd}}))$ is the command, with
servo parameters \texttt{--lookahead-time} ($0.1$) and \texttt{--servo-gain}
($300$).

\section{Backend realisation}\label{sec:backends}
Both backends receive the \emph{same} command from the loop above and only
differ in how they apply it.

\paragraph{RTDE backend (real robot / URSim).}
\begin{itemize}[leftmargin=1.5em]
  \item Speed: $\texttt{speedL}\big([\vv;\vw],\,a,\,2\Delta t\big)$.
  \item Servo: $\texttt{servoL}\big([\vp_{\text{cmd}};\rotvec(\quat{q}_{\text{cmd}})],\,0,0,\Delta t,\,t_{\text{look}},\,\text{gain}\big)$.
\end{itemize}

\paragraph{ROS backend (\texttt{launch\_all.sh} simulation, speed mode only).}
The Cartesian twist is mapped to joint velocities through the UR5 geometric
Jacobian $J(\bm{q})\in\R^{6\times6}$ (\texttt{kinematics.ur5\_fk\_and\_jacobian})
using damped least squares (\texttt{cartesian\_to\_joint\_vel}):
\begin{equation}
  \dot{\bm{q}} = J^{\top}\big(JJ^{\top} + \lambda^2 I\big)^{-1}
                 \begin{bmatrix}\vv\\\vw\end{bmatrix},
  \qquad \lambda=\texttt{damping}\ (0.05),
\end{equation}
followed by a joint-speed clamp to \texttt{max\_joint\_vel} ($1.5\,\mathrm{rad/s}$),
direction preserved. $\dot{\bm{q}}$ is published as a
\texttt{Float64MultiArray} on \texttt{/forward\_velocity\_controller/commands}.
Servo mode is unavailable here (no \texttt{servoL} equivalent) and is rejected
at the CLI.

\section{Safety}
\begin{itemize}[leftmargin=1.5em]
  \item \textbf{Grip deadman}: motion only while the grip button is held;
        releasing it re-baselines the origins on the next press.
  \item \textbf{Reader watchdog}: \texttt{oculus\_reader} serves cached poses
        with no timestamp, so staleness is detected by change-tracking the raw
        matrix. If the controller pose is unchanged for
        $>\texttt{watchdog\_timeout}$ ($0.25\,$s) while enabled, motion is
        stopped and a latch requires a deliberate grip re-press.
  \item \textbf{Stop ordering}: the mode-appropriate stop
        (\texttt{servoStop} in servo mode, else \texttt{speedStop}; zeros on
        the ROS backend) is issued on grip release, watchdog trip, and
        shutdown, \emph{before} \texttt{stopScript} and peripheral teardown.
        The ROS \texttt{forward\_velocity\_controller} holds the last command
        with no auto-decay, so publishing zeros is mandatory to stop.
\end{itemize}

\section{Symbol / flag / code reference}
\begin{center}\small
\begin{tabular}{@{}llll@{}}
\toprule
Symbol & Meaning & CLI flag & Default \\
\midrule
$s_{\text{pos}}$ & position scale & \texttt{--pos-scale} & $1.0$ \\
$s_{\text{prec}}$ & precision multiplier & \texttt{--precision-scale} & $0.5$ \\
$\alpha$ & filter constant & \texttt{--filter-alpha} & $0.8$ \\
$k_p,k_r$ & P-gains & \texttt{--pos-gain}, \texttt{--rot-gain} & $3.0,\ 2.0$ \\
$v_{\max},\omega_{\max}$ & velocity limits & \texttt{--max-lin-vel}, \texttt{--max-rot-vel} & $0.15,\ 0.75$ \\
$\delta_p,\delta_r$ & servo step clamps & \texttt{--max-pos-step}, \texttt{--max-rot-step} & $v_{\max}\Delta t,\ \omega_{\max}\Delta t$ \\
$t_{\text{look}}$ & servoL lookahead & \texttt{--lookahead-time} & $0.1$ \\
gain & servoL gain & \texttt{--servo-gain} & $300$ \\
$\lambda$ & IK damping (ROS) & \texttt{--damping} & $0.05$ \\
$M$ & axis remap & \texttt{--rmat-reorder} & $[-2,-1,-3,4]$ \\
\bottomrule
\end{tabular}
\end{center}

\end{document}
